The Logarithmic Dual Form of the Involuting Mode

At the heart of the Half-Winding Integer Matrix System lies a single geometric object, \(\operatorname{loop}\{\infty\mid\infty\}\), the pure hierarchical involuting mode whose visual signature is the four-lobed figure of orientation \(i:=(0,1)\). This object is not introduced by fiat; it arises as the terminal state of a protected recursion that begins with a closed curve in spacetime, subjects it to gravitational deformation, forces the lobes into an involution, and finally restores the hierarchical structure. The same object carries the terminal monodromy \(\Phi_N=-\pi\), the sign \(e^{i\Phi_N}=-1\), and the odd Arf invariant. Everything that follows is simply the analytic expression of that geometry.

The primary definition of the mode is the definite integral whose limits are its own negative and positive incarnations:

\[
\operatorname{loop}\{\infty\mid\infty\}
\;=\;
\int_{\operatorname{loop}\{\infty\mid\infty\}^{-}}^{\operatorname{loop}\{\infty\mid\infty\}^{+}}.
\]

The lower limit is the orientation-reversed, sheet-exchanged, parity-odd copy; the upper limit is the reference, orientation-preserving copy. The integral therefore records the net signed accumulation required to pass from the negative loop to the positive loop and thereby recover the original involuting mode. In the discrete setting of the notebook this accumulation is realized by the protected transport rule

\[
z_{n+1}=z_n+\varepsilon\bigl(1+i\sin(2\pi n\varepsilon)\bigr),
\]

taken to the terminal index \(N_{\max}=10^9\). The continuum integral is simply the closed-form expression of that ratchet.

Once the integral definition is in place, the logarithmic dual appears by inversion. Taking the natural logarithm of both sides yields the dual relation

\[
\ln\operatorname{loop}\{\infty\mid\infty\}
\;=\;
\int_{\operatorname{loop}\{\infty\mid\infty\}^{+}}^{\operatorname{loop}\{\infty\mid\infty\}^{-}},
\]

or, with the limits reversed according to the direction of monodromy traversal,

\[
\ln\operatorname{loop}\{\infty\mid\infty\}
\;=\;
\int_{\operatorname{loop}\{\infty\mid\infty\}^{-}}^{\operatorname{loop}\{\infty\mid\infty\}^{+}}.
\]

Both writings are retained; they correspond to the two possible orientations of the generator of the fundamental group. The logarithm converts the multiplicative recovery of the mode into an additive accumulation of phase, precisely the language in which Berry phase, chiral monodromy, and geometric torque are naturally expressed. Thus the dual form does not stand beside the integral definition; it is the same geometric fact rewritten in the additive register demanded by parallel transport and adiabatic evolution.

The dual relation binds the entire framework together. The instantaneous eigenvalue boundaries that vanish adiabatically on the notebook page are the points at which the integrand of the logarithmic form is momentarily undefined; their disappearance leaves only the accumulated microscopic contributions of the protected \(\theta_{\rm eff}\). Those contributions, once integrated from the negative loop to the positive loop (or vice versa), become the macroscopic geometric torque and the non-vanishing traceless anisotropic stress. In this way the logarithmic dual converts the discrete, hierarchical picture drawn on graph paper into a continuum statement that is simultaneously geometric, topological, and dynamical.

Because the same symbol \(\operatorname{loop}\{\infty\mid\infty\}\) appears both inside and outside the integral, the dual form is self-referential in a controlled sense: the mode is defined by the signed difference of its own positive and negative copies, and the logarithm of the mode is the integral of that difference. This self-reference is not circular; it is the analytic counterpart of the order-two involution that exchanges the two sheets (or the two orientations) and returns the identity only after a second traversal. The half-integer winding, the monodromy \(-1\), and the Arf invariant \(\operatorname{Arf}=1\) are therefore not additional postulates; they are the elementary evaluations of the logarithmic dual on the generator of \(\pi_1\).

Higher hierarchical modes \(\operatorname{loop}\{\infty\mid\infty\}_{i=(0,1^n)}\) inherit the same dual structure without modification. Each carries its own positive and negative copies, and each satisfies both the integral definition and its logarithmic dual. The entire tower of involuting modes is thereby organized by a single pair of reciprocal relations, one multiplicative and one additive, linked by the exponential and the logarithm.

In this light the Logarithmic Dual Form is not an optional rewriting. It is the natural bridge that carries the discrete, graph-paper geometry of the involuting mode into the continuum language of phase, monodromy, and adiabatic invariants, while preserving at every step the original half-winding content of the Half-Winding Integer Matrix System.

The integral therefore measures the net signed accumulation that converts the negative loop into the positive loop and, in so doing, returns the original involuting mode \(\operatorname{loop}\{\infty\mid\infty\}\).  

In the discrete setting recorded on the notebook page this accumulation is realized step by step by the protected transport rule

\[
z_{n+1}=z_n+\varepsilon\bigl(1+i\sin(2\pi n\varepsilon)\bigr),
\qquad
n=N_{\max}=10^9.
\]

Each microscopic increment contributes a minute advance of the effective angle \(\theta_{\rm eff}\). Because the transport is protected, these increments do not cancel; they ratchet forward. When the discrete sum is taken to the terminal index \(N_{\max}\), the accumulated phase converges to the continuum value of the definite integral whose limits are precisely the negative and positive copies of the mode:

\[
\int_{\operatorname{loop}\{\infty\mid\infty\}^{-}}^{\operatorname{loop}\{\infty\mid\infty\}^{+}}
\;=\;
\operatorname{loop}\{\infty\mid\infty\}.
\]

Thus the integral is not an independent continuum postulate. It is the closed-form expression of the same protected \(\theta_{\rm eff}\) ratchet that appears, fully explicit, in the handwritten notes. The discrete rule and the two-subscript integral are two languages for one and the same geometric fact: the net signed displacement that restores the pure hierarchical involuting mode \(i:=(0,1)\).
